\section{La plataforma OpenHands y la arquitectura del Agent}
\subsection{Componentes de la arquitectura de OpenHands}
OpenHands ofrece una plataforma de desarrollo de Agent contenedorizada y orquestable: el Agent puede operar simultáneamente la terminal, el editor y el navegador dentro de un entorno sandbox controlado; el backend puede conectarse a varios LLM (como code-davinci y modelos de código abierto), y a través de una cadena de herramientas de webhook e intérpretes de pruebas completa el ciclo cerrado de entrada→inferencia→acción.

Los componentes se comunican entre sí mediante un bus de eventos y una cola de mensajes, lo que permite sustituir el modelo o las herramientas en tiempo de ejecución. La plataforma también define metadatos para cada acción (forma en que se desencadena, contexto, salida esperada), lo que facilita la reversión y la depuración.

\subsection{Interfaz multimodal y auditoría}
\begin{figure}[H]
\centering
\includegraphics[width=\textwidth]{frame-01-openhands.jpg}
\caption{La pantalla de OpenHands muestra el flujo en cascada del Agent: percepción del entorno, inferencia del LLM y ejecución de herramientas. \protect\footnotemark}
\end{figure}
\footnotetext{Intervalo de tiempo en el video: 00:04:45--00:04:50.}

La plataforma pone énfasis en la «observabilidad»: cada acción de cada Agent (comandos de terminal, clics en el navegador, modificaciones de archivos) queda registrada, y los resultados de las pruebas y los registros fluyen de vuelta a un panel de auditoría central. OpenHands también ofrece interfaces para que una persona pueda interrumpir el proceso, usar la vista de depuración o insertar scripts auxiliares; esta capacidad de tipo plug-and-play hace que el Agent deje de ser una caja negra.

\begin{knowledgebox}{Capacidades centrales de OpenHands}
OpenHands incluye un adaptador multimodal, un motor de decisión basado en el LLM, ejecutores de herramientas (shell de Docker, pruebas, navegador) y un panel de auditoría visual. Cada parte se comunica mediante un bus de eventos, y el modelo o las herramientas pueden sustituirse en tiempo de ejecución.
\end{knowledgebox}

\begin{warningbox}{Observabilidad y límites de seguridad}
El Agent registra cada comando, pero también debe restringir las operaciones con privilegios elevados y garantizar que el sandbox pueda restaurarse; de lo contrario, los registros de desarrollo y el estado del entorno pueden filtrar información sensible y volverse difíciles de reproducir.
\end{warningbox}

\subsection{Resumen de la sección}
\begin{itemize}
    \item OpenHands conecta la percepción, la inferencia, la acción y la auditoría del Agent en un pipeline modular.
    \item Los registros transparentes y las interfaces de intervención humana son la garantía clave que hace que los equipos de ingeniería estén dispuestos a poner en producción un Agent.
\end{itemize}
